% !TEX program = xelatex
% Biên dịch: xelatex 04-experiment-plan-niche1.tex
% Tuân thủ 00-methodology.tex

\documentclass[11pt,a4paper]{article}

%=========================================================
% PACKAGES
%=========================================================
\usepackage[margin=2.2cm]{geometry}
\usepackage{fontspec}
\usepackage{polyglossia}
\setdefaultlanguage{vietnamese}
\setotherlanguage{english}
\setmainfont{Latin Modern Roman}
\setsansfont{Latin Modern Sans}
\setmonofont{Latin Modern Mono}

\usepackage{amsmath}
\usepackage{amssymb}
\usepackage{physics}
\usepackage{xcolor}
\usepackage{graphicx}
\usepackage{tabularx}
\usepackage{booktabs}
\usepackage{longtable}
\usepackage{array}
\usepackage{enumitem}
\usepackage{titlesec}
\usepackage{fancyhdr}
\usepackage{hyperref}
\usepackage{bookmark}
\usepackage{url}
\usepackage{listings}
\usepackage{tcolorbox}
\tcbuselibrary{skins, breakable, listings}

%=========================================================
% STYLE
%=========================================================
\definecolor{accent}{HTML}{0B5FA5}
\definecolor{accent2}{HTML}{B8860B}
\definecolor{success}{HTML}{2E8B57}
\definecolor{danger}{HTML}{B22222}
\definecolor{soft}{HTML}{F4F6FA}
\definecolor{softyellow}{HTML}{FFF8E1}
\definecolor{softgreen}{HTML}{E8F5E9}
\definecolor{codebg}{HTML}{F6F8FA}

\hypersetup{colorlinks=true, linkcolor=accent, urlcolor=accent, citecolor=accent2,
  pdftitle={04 - Experiment Plan Niche 1}, pdfauthor={TS. Do Phuc Hao}}

\titleformat{\section}{\Large\bfseries\color{accent}}{\thesection.}{0.5em}{}
\titleformat{\subsection}{\large\bfseries\color{accent}}{\thesubsection}{0.5em}{}
\titleformat{\subsubsection}{\normalsize\bfseries}{\thesubsubsection}{0.5em}{}

\setlist[itemize]{leftmargin=1.2em, itemsep=2pt, topsep=3pt}
\setlist[enumerate]{leftmargin=1.4em, itemsep=2pt, topsep=3pt}

\pagestyle{fancy}
\fancyhf{}
\rhead{\small\textit{Experiment Plan --- Niche 1}}
\lhead{\small TS. Đỗ Phúc Hảo}
\rfoot{\thepage}
\renewcommand{\headrulewidth}{0.3pt}

\newtcolorbox{note}[1][]{enhanced, breakable, colback=soft, colframe=accent,
  boxrule=0.4pt, arc=2pt, left=6pt, right=6pt, top=4pt, bottom=4pt,
  title={#1}, fonttitle=\bfseries\color{white}, coltitle=white, colbacktitle=accent}

\newtcolorbox{greenbox}[1][]{enhanced, breakable, colback=softgreen, colframe=success,
  boxrule=0.6pt, arc=2pt, left=6pt, right=6pt, top=4pt, bottom=4pt,
  title={#1}, fonttitle=\bfseries\color{white}, coltitle=white, colbacktitle=success}

\newtcolorbox{goldbox}[1][]{enhanced, breakable, colback=softyellow, colframe=accent2,
  boxrule=0.6pt, arc=2pt, left=6pt, right=6pt, top=4pt, bottom=4pt,
  title={#1}, fonttitle=\bfseries\color{white}, coltitle=white, colbacktitle=accent2}

\lstdefinestyle{codestyle}{backgroundcolor=\color{codebg},
  basicstyle=\ttfamily\small, breaklines=true, frame=single, framerule=0.3pt,
  rulecolor=\color{gray!40}, xleftmargin=0.5em, xrightmargin=0.5em,
  aboveskip=4pt, belowskip=4pt, numbers=none, tabsize=2, showstringspaces=false}
\lstset{style=codestyle}

%=========================================================
% METADATA
%=========================================================
\title{\vspace{-1.5cm}\textbf{\color{accent} 04 --- Experiment Plan} \\[3pt]
       \Large DynaQGNN-IDS cho NTN / 6G SAGIN \\[2pt]
       \normalsize \textit{Dataset, Baselines, Metrics, Ablation, Compute Budget}}
\author{TS. Đỗ Phúc Hảo \\ \small Khoa CNTT --- Đại học Kiến trúc Đà Nẵng}
\date{Ngày 24 tháng 04 năm 2026 (phiên bản 0.1)}

\begin{document}
\maketitle
\thispagestyle{fancy}

\begin{note}[Mục đích]
Tài liệu cụ thể hoá \textbf{kế hoạch thực nghiệm} cho paper A (warm-up, WCL/OJ-COMS) và paper B (full, TWC/T-IFS). Bao phủ: dataset + simulator, baseline suite, metric, ablation matrix, compute budget (RTX 4090 + Colab Pro+), lộ trình chạy theo tuần. Tài liệu \textit{ràng buộc} đối với code repo sẽ tạo ở bước 5.
\end{note}

\tableofcontents
\vspace{0.4cm}
\hrule
\vspace{0.4cm}

%=========================================================
\section{Phạm vi thí nghiệm của paper A vs paper B}
%=========================================================

\begin{longtable}{@{}p{3.5cm} p{5.5cm} p{5.5cm}@{}}
\toprule
\textbf{Mục} & \textbf{Paper A (warm-up)} & \textbf{Paper B (full)} \\
\midrule
\endhead
Target & IEEE WCL / OJ-COMS & IEEE T-IFS / TWC / IoTJ \\
Đóng góp & Đóng góp 1 (DynaQGNN centralized) & Cả 3 đóng góp + federated + benchmark \\
Scale satellite $K$ & $\leq 20$ & $20 \rightarrow 100 \rightarrow 500$ (scaling exp) \\
Dataset & 1 (CICIDS2017 + NTN mapping) & 3 (CICIDS, ToN-IoT, Bot-IoT) $+$ Hypatia-DoS \\
Baseline & 4 (Random, E-GraphSAGE, QSVC, QGNN-vanilla) & 8 (thêm Anomal-E, E-GRACL, GAT, QGAN-IDS) \\
Ablation & 3 (qubits, depth, encoding) & 6 (thêm ansatz, noise, FL config) \\
Qubits max & 16 & 24--28 \\
Số seed & 3 (cho draft), 5 (cho submit) & 5 bắt buộc \\
Wall-clock budget & $\leq$ 2 tuần RTX 4090 & $\leq$ 6 tuần RTX 4090 + Colab \\
\bottomrule
\end{longtable}

%=========================================================
\section{Dataset \& simulator}
%=========================================================

\subsection{Tầng 1 --- Flow dataset (nội dung traffic)}

\begin{longtable}{@{}p{3cm} p{3cm} p{9cm}@{}}
\toprule
\textbf{Dataset} & \textbf{Size} & \textbf{Ghi chú sử dụng} \\
\midrule
\endhead
CICIDS-2017 & $\sim$3M flows & Có 14 attack types. \textbf{Dataset chính paper A}. Preprocessing: Zeek/NetFlow features. \\
ToN-IoT & $\sim$22M records & IoT-IIoT, có telemetry + network. \textbf{Dùng paper B} cho IoT angle. \\
Bot-IoT & $\sim$73M records & Focus botnet + DDoS. \textbf{Dùng paper B} cho stealthy DDoS. \\
CSE-CIC-IDS2018 & $\sim$16M & Thay thế tuỳ chọn CICIDS-2017 nếu reviewer yêu cầu dataset mới hơn. \\
\bottomrule
\end{longtable}

\subsection{Tầng 2 --- NTN simulator (topology \& attack injection)}

\begin{longtable}{@{}p{3cm} p{12cm}@{}}
\toprule
\textbf{Tool} & \textbf{Vai trò} \\
\midrule
\endhead
\textbf{Hypatia} & Base simulator (Rus et al.). Chạy ns-3 packet-level + quỹ đạo Kepler. \textbf{Constellation chuẩn}: Starlink shell-1 (72 plane $\times$ 22 sat = 1584 sat ở 550 km). Paper A chỉ dùng sub-constellation $K \leq 20$. \\
\textbf{Hypatia-DoS ext. 2025} & Plugin mở rộng 4 attack: TCP SYN flood, UDP flood, ICMP flood, Shrew low-rate pulsing. \textbf{Dùng nguyên}. \\
\textbf{xeoverse} & Dự phòng nếu cần scale $K > 500$ với real-time --- nhanh 2.9--40$\times$. \\
\bottomrule
\end{longtable}

\subsection{Pipeline dataset hoá}

\begin{enumerate}
  \item Chạy Hypatia + Hypatia-DoS trên constellation thu nhỏ ($K=20$) với 30 phút mô phỏng.
  \item Xuất \textbf{packet trace} + \textbf{satellite-link topology theo thời gian}.
  \item Gắn \textbf{flow-level features} từ CICIDS-2017 (resp. ToN-IoT) vào các \textit{luồng logic} giữa ground station $\rightarrow$ satellite $\rightarrow$ ground gateway.
  \item Xây dựng \textbf{graph snapshot} mỗi $\Delta t = 1$ s: nodes = $\{$GS, sat, gateway$\}$, edges = ISL + uplink + downlink.
  \item Nhãn: mỗi flow có nhãn (benign / attack type) kế thừa từ CICIDS-2017.
  \item Chia train/val/test theo thời gian \textbf{không random} (time-split để tránh leakage).
\end{enumerate}

\begin{goldbox}[Quyết định quan trọng về graph construction]
Graph của ta là \textbf{heterogeneous, temporal}:
\begin{itemize}
  \item Node types: $\{$UE, LEO-sat, gateway$\}$.
  \item Edge types: $\{$uplink, ISL, downlink$\}$ --- có feature riêng (throughput, delay, SNR).
  \item Adjacency $\mathbf{A}_t$ thay đổi theo orbital dynamics $\Rightarrow$ dùng \textbf{temporal snapshot sequence}.
\end{itemize}
Đây là điểm khác biệt so với E-GraphSAGE (static, homogeneous) $\Rightarrow$ phải làm nổi bật trong paper.
\end{goldbox}

%=========================================================
\section{Baseline suite}
%=========================================================

\subsection{Nhóm 1 --- Trivial baselines}
\begin{itemize}
  \item Random classifier --- sanity check.
  \item Logistic regression trên flow features --- lower bound.
  \item MLP (2--3 layers) trên flow features --- tabular baseline.
\end{itemize}

\subsection{Nhóm 2 --- Classical GNN baselines}
\begin{itemize}
  \item \textbf{E-GraphSAGE} (Lo et al. 2021) --- SOTA graph IDS. Code: \url{https://github.com/waimorris/E-GraphSAGE}.
  \item \textbf{Anomal-E} (Caville et al. 2022) --- self-supervised baseline.
  \item \textbf{E-GRACL} (2024) --- contrastive GNN.
  \item \textbf{GAT-IDS} (Nature SR 2025) --- attention mechanism.
  \item \textbf{TGN} (Temporal Graph Networks, Rossi 2020) --- temporal baseline bắt buộc vì ta có temporal graph.
\end{itemize}

\subsection{Nhóm 3 --- Quantum baselines}
\begin{itemize}
  \item \textbf{QSVC / QuIDS} (4 qubits, 8 features) --- kernel method baseline.
  \item \textbf{QGAN-IDS} (2025, NSL-KDD style) --- generative baseline.
  \item \textbf{QGNN-vanilla} (Recy 2026-style) --- same QGNN but without temporal encoding.
  \item \textbf{Quantum AE} (QMI 2024) --- unsupervised baseline.
\end{itemize}

\subsection{Nhóm 4 --- Proposed (ablation nội bộ)}
\begin{itemize}
  \item \textbf{DynaQGNN-IDS} (full) --- proposed.
  \item \textbf{DynaQGNN without temporal encoding}.
  \item \textbf{DynaQGNN without barren-plateau mitigation}.
  \item \textbf{DynaQGNN with random ansatz (không problem-inspired)}.
\end{itemize}

%=========================================================
\section{Metric}
%=========================================================

\subsection{Detection metrics}
\begin{itemize}
  \item \textbf{Binary}: Accuracy, Precision, Recall, F1, AUC-ROC, AUC-PR (ưu tiên PR vì lớp attack hiếm).
  \item \textbf{Multiclass}: Macro-F1 + per-class F1 (đặc biệt cho Shrew low-rate).
  \item \textbf{Miss / False Alarm}: $\mathbb{P}[\mathrm{miss}], \mathbb{P}[\mathrm{FA}]$ theo ngưỡng $\tau$.
\end{itemize}

\subsection{Efficiency metrics}
\begin{itemize}
  \item Wall-clock training time (s / epoch).
  \item Inference latency (ms) trên 1 graph snapshot.
  \item Peak GPU memory (GB).
  \item Circuit depth + qubits + CNOT count.
\end{itemize}

\subsection{Robustness metrics}
\begin{itemize}
  \item Performance dưới CSI error $\sigma_e \in \{0.01, 0.05, 0.1, 0.2\}$.
  \item Performance dưới attack burstiness $\pi_t$ khác nhau.
  \item Performance khi $K$ thay đổi $\{5, 10, 20, 50, 100\}$ --- scalability.
  \item Performance khi topology mutation rate tăng (satellite handover frequency).
\end{itemize}

\subsection{Quantum-specific diagnostics}
\begin{itemize}
  \item \textbf{Barren plateau}: variance of gradient $\mathrm{Var}(\partial_\theta \mathcal{L})$ theo số qubit.
  \item \textbf{Expressibility}: KL divergence giữa distribution của circuit và Haar-random (Sim et al. 2019).
  \item \textbf{Entanglement capability}: Meyer-Wallach measure.
  \item \textbf{Noise resilience}: performance dưới depolarizing noise $p \in \{0, 0.001, 0.005, 0.01\}$.
\end{itemize}

%=========================================================
\section{Ablation matrix}
%=========================================================

\begin{longtable}{@{}p{0.6cm} p{3.8cm} p{10.6cm}@{}}
\toprule
\textbf{\#} & \textbf{Yếu tố} & \textbf{Các giá trị / cấu hình} \\
\midrule
\endhead
A1 & Số qubits & $\{4, 8, 12, 16, 20, 24\}$ --- paper A dừng ở 16, paper B đến 24 \\
A2 & Circuit depth (variational layers) & $\{2, 4, 6, 8\}$ \\
A3 & Encoding scheme & \{angle, amplitude, re-uploading\} \\
A4 & Ansatz family & \{hardware-efficient, problem-inspired (topology-aware)\} \\
A5 & Barren plateau mitigation & \{none, layer-wise training, local cost, Gaussian init\} \\
A6 & Temporal encoding & \{disabled, 1-step, 3-step, RNN-style\} \\
A7 & Noise model (paper B) & \{noiseless, depolarizing $p$, IBM realistic backend\} \\
A8 & FL configuration (paper B) & \{centralized, FedAvg, FedProx, SCA-weighted\} \\
\bottomrule
\end{longtable}

\textbf{Tổng số run ước tính paper A}: $(3 \text{ ablations chính} \times 5 \text{ values each}) \times 3 \text{ seeds} \approx 45$ runs. Mỗi run $\sim 30$ min on RTX 4090 $\Rightarrow$ $\approx 22$ GPU-hours. Khả thi trong 1 tuần.

\textbf{Tổng số run ước tính paper B}: $\approx 250$ runs $\times 20$ min trung bình $\approx 83$ GPU-hours. Khả thi trong 4--5 tuần RTX 4090 + sweep trên Colab Pro+ song song.

%=========================================================
\section{Compute budget chi tiết}
%=========================================================

\subsection{RTX 4090 (máy chính)}
\begin{itemize}
  \item 24 GB VRAM, Ada Lovelace, cuQuantum compatible.
  \item Chạy \texttt{pennylane-lightning-gpu}: simulator 24 qubits $\sim$5--10 s / forward pass (với 4--6 variational layers).
  \item Training 1 epoch trên 100k graph snapshots: $\sim$15--20 min với 12 qubits.
  \item Checkpoint mỗi 5 epoch, tổng epoch 30--50.
  \item \textbf{Dùng cho}: Ablation chính, baseline chính, reproducibility cuối.
\end{itemize}

\subsection{Colab Pro+ (dự phòng)}
\begin{itemize}
  \item A100 40 GB khi may mắn (hoặc V100 16 GB).
  \item \textbf{Dùng cho}: Hyperparameter sweep song song, quick sanity check khi không có access RTX 4090.
  \item \textbf{Tránh dùng cho}: final numbers --- instability về GPU allocation.
\end{itemize}

\subsection{IBM Quantum Experience}
\begin{itemize}
  \item Free tier: 7 qubits Nairobi, Lagos; 127 qubits Sherbrooke (queue).
  \item \textbf{Chỉ dùng}: Validation final trên 1--2 configuration nhỏ (4--8 qubits) để có ``real hardware'' bullet cho paper B.
  \item \textbf{Không dùng cho}: training.
\end{itemize}

\subsection{Storage}
\begin{itemize}
  \item Dataset raw: $\sim$60 GB (CICIDS + ToN-IoT + Bot-IoT + Hypatia traces).
  \item Checkpoints: $\sim$5 GB / 50 checkpoints.
  \item W\&B artifacts: online.
\end{itemize}

%=========================================================
\section{Lộ trình theo tuần (12 tuần cho paper A)}
%=========================================================

\begin{longtable}{@{}p{1.3cm} p{14cm}@{}}
\toprule
\textbf{Tuần} & \textbf{Nhiệm vụ} \\
\midrule
\endhead
1 & Setup RTX 4090 env (conda, pennylane-lightning-gpu, cuQuantum). Tạo repo GitHub. \\
2 & Chạy Hypatia + Hypatia-DoS; xuất sample dataset 1 hour constellation $K=10$. Sanity-check loader. \\
3 & Implement classical baselines: Logistic, MLP, E-GraphSAGE. Chạy baseline 1 seed. \\
4 & Implement Anomal-E, TGN (temporal). Chạy baseline 3 seeds. \\
5 & Implement DynaQGNN v0 (centralized, 8 qubits, no barren mitigation). Chạy pilot. \\
6 & Chạy ablation A1 (qubits 4--16) + A2 (depth 2--8). Phân tích kết quả. \\
7 & Ablation A3 (encoding) + A5 (barren mitigation). Chọn cấu hình best. \\
8 & Chạy main table với 3 seed, 5 seed cho top-3 cấu hình. \\
9 & Draft paper v0.1: Abstract, Intro, Related Work, System Model, Problem Formulation. \\
10 & Draft paper v0.2: Proposed Method, Experiments, Results. \\
11 & Hoàn thiện figures (TikZ cho system model, convergence, ablation). Proofreading. \\
12 & Iteration cuối + submit WCL / OJ-COMS. \\
\bottomrule
\end{longtable}

Paper B bắt đầu ngay tuần 8 song song (federated + scaling experiment), dự kiến submit tháng 9--12.

%=========================================================
\section{Reproducibility checklist (bắt buộc theo 00-methodology)}
%=========================================================

\begin{itemize}
  \item[$\square$] Random seed cố định trong config YAML.
  \item[$\square$] 5 seed khác nhau cho số liệu công bố.
  \item[$\square$] Git commit hash log vào W\&B mỗi run.
  \item[$\square$] Config YAML full, no magic number trong code.
  \item[$\square$] Dataset hash / DVC version.
  \item[$\square$] Training time + GPU log.
  \item[$\square$] Environment.yml pinned cho conda + requirements.txt cho pip.
  \item[$\square$] Colab bootstrap notebook test được trên máy trống.
  \item[$\square$] Code release link trong paper (GitHub public sau accept, anonymous github trong review).
  \item[$\square$] Numerical results (CSV) release kèm figures để reviewer chạy lại plot.
\end{itemize}

%=========================================================
\section{Rủi ro thực nghiệm và cách giảm}
%=========================================================

\begin{longtable}{@{}p{4.5cm} p{10.5cm}@{}}
\toprule
\textbf{Rủi ro} & \textbf{Cách giảm} \\
\midrule
\endhead
QGNN không đánh bại classical GNN & Điều chỉnh regime: tăng topology mutation rate (nơi classical dễ fail); đo trên extreme non-IID / few-shot; tập trung vào \textit{parameter efficiency} thay vì raw accuracy. \\
Barren plateau vô hiệu hoá training & Apply đầy đủ Cerezo 2021 + Grant 2019; dùng shallow ansatz (depth $\leq 4$); layer-wise training. \\
Hypatia-DoS không có đủ attack variety & Bổ sung synthetic attack bằng GAN / rule-based trên dataset flow; hoặc inject malware flow từ ToN-IoT. \\
Compute blow-up ở ablation A7 noise & Giới hạn noise study ở 12 qubits, 4 layers; ngoại suy phần còn lại. \\
Reviewer đòi real-QPU validation & Dành 1 sub-experiment chạy IBM Nairobi 7 qubits / Sherbrooke --- chỉ cần 1 configuration. \\
Thiếu time cho paper A (trễ timeline) & Scope paper A chỉ centralized + 1 dataset; paper B mở rộng; không trộn. \\
\bottomrule
\end{longtable}

%=========================================================
\section{Cấu trúc repo sẽ tạo ở bước 5}
%=========================================================

(Tham chiếu \texttt{00-methodology.tex} Section 7; đây là bản cụ thể hoá cho niche này)

\begin{lstlisting}[language=bash]
dynaqgnn-ntn-ids/
|-- README.md
|-- environment.yml
|-- requirements.txt
|-- Makefile
|-- .gitignore
|-- src/
|   |-- data/
|   |   |-- hypatia_pipeline.py     # chay Hypatia + extract
|   |   |-- graph_builder.py        # build temporal graph
|   |   `-- flow_mapper.py          # map CICIDS flows vao topology
|   |-- models/
|   |   |-- baselines_classical.py  # MLP, E-GraphSAGE, Anomal-E, TGN
|   |   |-- baselines_quantum.py    # QSVC, QGAN, Quantum AE
|   |   |-- dynaqgnn.py             # proposed
|   |   `-- ansatz_library.py       # collection of ansatze
|   |-- training/
|   |   |-- trainer.py
|   |   |-- callbacks.py
|   |   `-- bp_mitigation.py        # barren plateau utilities
|   `-- utils/
|       |-- metrics.py              # all metrics in Section 5
|       |-- expressibility.py
|       `-- seeds.py
|-- experiments/
|   |-- exp01_classical_baselines/
|   |-- exp02_quantum_baselines/
|   |-- exp03_dynaqgnn_centralized/
|   |-- exp04_ablation_qubits/
|   |-- exp05_ablation_depth/
|   |-- exp06_ablation_encoding/
|   |-- exp07_ablation_bp/
|   `-- exp08_main_table/
|-- configs/
|   |-- base.yaml
|   |-- baselines/
|   `-- dynaqgnn/
|-- notebooks/
|   |-- 01_explore_hypatia_output.ipynb
|   |-- 02_graph_construction_demo.ipynb
|   `-- colab_bootstrap.ipynb
|-- data/        # gitignored
|-- results/     # gitignored, W&B fallback
`-- paper/
    |-- main.tex
    |-- refs.bib
    `-- figures/
\end{lstlisting}

%=========================================================
\section{Hành động ngay sau tài liệu này}
%=========================================================

\begin{enumerate}
  \item \textbf{File 05}: \texttt{05-tier-a-reading-note.tex} --- tóm tắt 4 foundation paper (Biamonte 2017, Schuld 2021, Cerezo 2021, McClean 2018). Thời gian: $\approx 1$ phiên.
  \item \textbf{Figures 1 + 2}: \texttt{figures/fig-system-model.tex} + \texttt{figures/fig-dynaqgnn-ansatz.tex} bằng TikZ (theo Gem 2).
  \item \textbf{Repo skeleton}: \texttt{dynaqgnn-ntn-ids/} với cấu trúc trên, \texttt{environment.yml}, \texttt{Makefile}, \texttt{colab\_bootstrap.ipynb}, và \texttt{src/data/hypatia\_pipeline.py} ở mức stub.
  \item \textbf{Push commit đầu tiên} lên GitHub repo (private ban đầu, public sau khi submit).
\end{enumerate}

\begin{greenbox}[Sẵn sàng bắt đầu code]
Tới đây, chúng ta có đủ để chuyển từ ``giấy'' sang ``code'': mục tiêu rõ (3 đóng góp), dataset chọn (Hypatia + CICIDS), baseline liệt kê (13 biến thể), metric định nghĩa, compute budget biết, timeline 12 tuần.
\end{greenbox}

\vspace{0.4cm}
\hrule
\vspace{0.3cm}
\textit{\small Phiên bản 0.1 --- 24/04/2026. Cập nhật khi thay đổi scope hoặc compute budget.}

\end{document}
